% Generated draft: requires mathematical and source review.
\documentclass[11pt]{article}
\usepackage[margin=25mm]{geometry}
\usepackage{fontspec}
\setmainfont{Latin Modern Roman}
\usepackage{amsmath,amssymb,mathtools,unicode-math}
\setmathfont{Latin Modern Math}
\usepackage{enumitem,needspace,xurl,hyperref}
\setlength{\parindent}{0pt}
\setlength{\parskip}{5pt}
\setlength{\emergencystretch}{3em}
\newcommand{\N}{\mathbb{N}}
\newcommand{\Z}{\mathbb{Z}}
\newcommand{\Q}{\mathbb{Q}}
\newcommand{\R}{\mathbb{R}}
\newcommand{\C}{\mathbb{C}}
\newcommand{\F}{\mathbb{F}}
\newcommand{\A}{\mathbb{A}}
\newcommand{\PP}{\mathbb{P}}
\newcommand{\E}{\mathbb{E}}
\newcommand{\eps}{\varepsilon}
\newcommand{\dd}{\,\mathrm{d}}
\begin{document}
\section*{UnsolvedMath ID: 4000002}
\textbf{Source problem code:} catalogue entry
\subsection{Definitions and mathematical statement}
For \(N\ge1\) and \(1\le p\le\infty\), let \(d_p(x,y)=\|x-y\|_p\) on \(\R^N\), with the usual maximum norm when \(p=\infty\). Coarse Ricci curvature requires not only a metric but a family of probability measures \((m_x)\), for example the transition laws of a random walk. For measures of finite first moment, set
\[
W_1(\mu,\nu)=\inf_{\pi\in\Pi(\mu,\nu)}\int d(u,v)\,\mathrm d\pi(u,v),
\qquad
\kappa(x,y)=1-\frac{W_1(m_x,m_y)}{d(x,y)}\quad(x\ne y),
\]
where \(\Pi(\mu,\nu)\) is the set of couplings. Thus the flat observation \(\kappa(x,y)=0\) means \(W_1(m_x,m_y)=d_p(x,y)\). A customary flat-space choice is a fixed finite-first-moment law translated with the base point, \(m_x=(z\mapsto x+z)_*\nu\); the extracted source does not state its choice of laws or scale.

Ollivier's Problem B asks whether this flat \(L^p\)-norm observation yields useful results for Finsler manifolds, whose tangent spaces carry varying Finsler norms or gauges and whose distance is obtained by integrating them along curves. This is a research programme, not a specified universal theorem: the record does not identify a desired inequality, invariant, classification, or other conclusion. Nor does it prescribe a reference measure, transition laws and scale on a Finsler manifold, regularity/convexity assumptions, or treatment of a possibly nonreversible distance. Consequently a unique precise yes/no assertion beyond the stated programme cannot be extracted without adding mathematical goals.

The supplied research summary reports partial progress through displacement-convexity approaches of Ohta and Ohta--Sturm, while distinguishing that work from a comprehensive theory based on this specific coarse-Ricci observation. It supplies no theorem asserting that the requested programme has been completed.
\subsection{Short English statement}
Can the vanishing of coarse Ricci curvature in normed \(\R^N\) lead to useful theorems for Finsler manifolds? The source leaves the Finsler random walk and the desired mathematical conclusions unspecified.
\subsection{Sources}
\begin{itemize}
\item Local UnsolvedMath registry, permanent ID \(4000002\), exact source code \texttt{catalogue entry}, ``Finsler manifolds.''
\item Yann Ollivier, \emph{Discrete Ricci curvature: Open problems} (2008), Problem B; \url{https://www.yann-ollivier.org/rech/publs/problems_curvmarkov.pdf}. Recorded access 2026-07-29; pdf-text extraction.
\item Supplied research summary mentions Ohta and Ohta--Sturm on displacement convexity in Finsler manifolds. It provides no complete publication references or precise results.
\end{itemize}
\end{document}
